\section{Introduction}
\textit{``What magical trick makes us intelligent? The trick is that there is no trick. The power of intelligence stems from our vast diversity, not from any single, perfect principle.''}

\qquad\qquad\qquad\qquad\qquad\qquad\qquad\qquad\qquad\qquad\qquad \textit{- Marvin Minsky, The Society of Mind, p. 308}

Confronted with the complexities of real-world tasks, solving them often requires multiple steps. The rapid progress of chat-based large-scale language models (LLMs) has yielded remarkable achievements in complex task-solving \cite{openai_chatgpt,ouyang2022training,thoppilan2022lamda,pichai_2023,claude,bai2022constitutional,wei2022emergent,bubeck2023sparks}. Nevertheless, it is worth noting that their success is heavily reliant on human input to guide the conversation in the right direction. This reliance necessitates users to provide relevant and precise prompts based on their intentions and the chat agent's feedback. This can be challenging, time-consuming, and sometimes impossible. Crafting effective prompts often demands a deep understanding and expertise of a particular domain of knowledge.
Consider an individual who lacks trading expertise; they would find it difficult to create suitable prompts for directing a 
chat agent to develop a trading application. This predicament is raising a crucial question: can we replace human intervention with an autonomous communicative agent capable of steering the conversation toward task completion with minimal
human supervision? To tackle this issue, it is crucial to conduct more research exploring the potential, capabilities, and limitations of communicative agents that operate entirely on their own to complete tasks. Understanding how multiple agents interact with each other is important for anticipating the future of artificial intelligence.
The dynamics of collaborating or competing agents play a key role in determining the success of AI systems \cite{asimov1940robot,dafoe2021cooperative,dafoe2020open,ouyang2022training,saunders2022self,bai2022training,bai2022constitutional}.

This paper explores the potential of building scalable techniques to facilitate autonomous cooperation among communicative agents and provide insight into their ``cognitive'' processes. Several challenges arise when asking a society of agents to autonomously cooperate on completing tasks. Examples we encountered in our preliminary analysis include \textit{role flipping}, \textit{assistant repeating instructions}, \textit{flake replies}, and \textit{infinite loop of messages}.
Therefore, it is critical to investigate ways to align these models with human intentions and to explore means enabling their effective cooperation.
To address these issues, we propose a novel cooperative agent framework named \RP to automate cooperation between communicative agents. Specifically, our proposed approach involves using \RP with \textit{inception prompting} to autonomously guide the communicative agents toward task completion. 
Only a preliminary \textit{idea} is needed from human 
to guide the conversations toward complex task-solving.

Our library, which we make publicly available, provides modular functionality, and includes implementations of different agents, examples of well-crafted prompts, and data explorers.
We hope our library serves as a ground for future research in various areas such as multi-agent systems, cooperative AI, game theory simulations, social analysis, AI ethics, AI alignment, and beyond.

In addition, our \RP method provides a highly scalable way to generate conversational data for studying the behaviors and capabilities of chat agents. We showcase how \RP can be used to let chat agents communicate with each other for task completion and record their conversations for behavior analysis and capability understanding. 
In particular, we consider two cooperative scenarios of role-playing and generate two large conversational, task-oriented, and instruction-following datasets: \textit{AI Society} and \textit{Code}. We also use our framework to collect two single-turn question-answer datasets, \textit{Math} and \textit{Science}, for LLM ability emergence study. Furthermore, we generate a \textit{Misalignment} dataset that is a simulation of possible malicious applications which demonstrate the potential risks of an unaligned autonomous agent system. 
The datasets offer a valuable resource for investigating conversational language models, enabling them to comprehend and react to human language more effectively. Furthermore, our \RP offers a scalable method of creating conversational instruction-following data, which can potentially enhance the development of more advanced 
language models. We show that solutions derived from our \RP framework outperform those generated in a single shot by \texttt{gpt-3.5-turbo} \cite{openai_chatgpt} in both GPT4 and human evaluations. We also study knowledge emergence in LLMs by fine-tuning LLaMA \cite{touvron2023llama} on progressively growing datasets generated through our framework. Additionally, we evaluate our code generation capabilities through benchmarking our final model on HumanEval \cite{chen2021evaluating} and HumanEval$^{+}$ \cite{evalplus}. 


\mysection{Contributions} Our contributions are fourfold:
(1) We introduce a novel cooperative agent framework, \RP, that allows communicative agents to collaborate autonomously toward completing tasks while requiring minimal human intervention;
(2) Our framework offers a scalable approach for studying the cooperative behaviors and capabilities of multi-agent systems. It illuminates the challenges of achieving autonomous cooperation, and provides strategies for addressing them. We showcase the potential power of multi-agent collaboration for complex-task solving;
(3) We demonstrate the significant emergence of LLM training abilities by utilizing the datasets we have collected from simulating four distinct agent collaboration scenarios;
(4) We have open-sourced our library, containing implementations of various agents, data generation pipelines, data analysis tools, and collected datasets, to support research on communicative agents and beyond.
